% OPERA manuscript body, shared by the article (opera_manuscript.tex) and the
% thesis (Paper III). Edit here; both pick it up.
\begin{abstract}
Clinical prediction models are usually built one task at a time, for one
disease, one outcome and one dataset. Each learns only from its own labelled patients, so
most questions have too few events for a model of their own, even in common
diseases, and the signal that related tasks hold goes unused. Multi-task learning
and pretraining can share that signal, but it is rarely known where the shared
signal comes from, how much each source adds, or whether a model can learn which
sources to trust. We tested this across a whole specialty, where tasks are many
and each is small. In a nationwide cohort of 37,531 Danish adults with 24
lymphoid cancer subtypes, sequence models read pre-treatment register histories
and predicted 69 outcomes, from treatment failure to infection, organ toxicity
and acute care, at horizons from 30 days to two years. A matched design widened
the sources of signal step by step: other diseases, other outcomes, both, a
model that learns what to borrow from each, and an encoder pretrained on the
records of the whole Danish population. With every decision fixed on patients
treated until 2022, the models were evaluated once, on 6,446 patients treated
from 2023, against gradient-boosted trees tuned for every task and the
prognostic indices in clinical use. \opera{}, the model that performed best on
the tuning year, was more accurate than the tuned trees in
\result{T2_better}{X of Y} disease--outcome predictions (median AUC difference
\result{T2_median}{Z}) and did the work of \result{narrow_fits}{N} separately
trained models. \placeholder{Other diseases and other outcomes each added
signal, learning what to borrow added X beyond sharing everything, and the gains
were largest where a prediction had least data of its own; population
pretraining changed performance by Y.} In diffuse large B-cell lymphoma, at the
specificity of the IPI, \opera{} identified \result{ipi_extra_per1000}{N} more
per 1,000 of the patients who died within two years and placed
\result{ipi_low_extra_per1000}{M} more at low risk. Its predictions rested on
findings clinicians recognise, read differently for each complication.
\placeholder{Draft if the results hold: clinical prediction need not be built
one task at a time, and letting a model learn what to borrow can extend it to
questions too small for a model of their own.}
\end{abstract}

\noindent\fbox{\parbox{\dimexpr\textwidth-2\fboxsep-2\fboxrule}{%
\small\textbf{Key findings} \placeholder{(confirm each against the locked test-set analysis)}
\begin{itemize}[leftmargin=1.2em,itemsep=1pt,topsep=3pt]
  \item One model predicted 69 outcomes in 24 lymphoid cancer subtypes and was
  more accurate than per-task XGBoost in \result{T2_better}{X of Y}
  disease--outcome predictions, covering the tasks of
  \result{narrow_fits}{N} separately trained models.
  \item \placeholder{Learning across diseases and across outcomes each improved
  prediction, the two added up, letting the model learn what to borrow added
  more, and the gains were largest where data were scarcest.}
  \item In DLBCL, at the specificity of the IPI, it identified
  \result{ipi_extra_per1000}{N} more per 1,000 of the patients who died within
  two years and placed \result{ipi_low_extra_per1000}{M} more at low risk with
  an observed risk no higher than the IPI low-risk group.
  \item \placeholder{Nation-scale pretraining improved specialised prediction
  directly / after adaptation to the cohort, and reached the full-data
  cohort-pretrained model within the prespecified margin using X\% of its
  labelled training patients.}
  \item It relied on findings clinicians recognise, read differently for each
  complication, and can explain each patient's risks by the episodes of care
  behind them.
  \item All results are from one evaluation of patients treated from 2023, after
  every modelling decision had been fixed and while first-line treatment was
  changing.
\end{itemize}}}

\operasection{Introduction}
Clinical prediction models are usually built one task at a time. Each is developed for
one disease and one outcome, from the patients who have both recorded, and the
next question starts again from scratch. A model learned this way can only be
as good as the events behind it allow \citep{Riley2020}, and events are scarce
even where patients are not: most outcomes are uncommon, and flexible models
need many more events than regression to perform well
\citep{vanderploeg2014data}. Prediction improves steadily as data grow
\citep{hestness2017scaling}, yet every task-specific model leaves most of the
available data unused. Clinical tasks are not independent. Patients with
different diseases share physiology, treatments, complications and patterns of
care, and different outcomes in the same patient are shaped by the same frailty,
prior illness and course of treatment. Much of what a prediction model must
learn is how to read a patient's record, and there is reason to expect that
reading to change less from task to task than the target does. Learned once
from every task, it would leave each task to learn mainly what is particular to
it \citep{Baxter2000,Maurer2016}, much
as partial pooling trusts each estimate in proportion to its evidence
\citep{EfronMorris1975}. This is the premise of multi-task learning and
pretraining, and it suggests a different default for clinical prediction:
share signal across tasks, and let the model learn what to borrow from where.
Like players in an orchestra, each prediction is then made in the company of
the others rather than in isolation.

The idea itself is well established. Multi-task learning trains related
predictions together
\citep{caruana1997multitask}, and self-supervised pretraining learns from
recorded histories before any outcome is labelled, as BEHRT, Med-BERT,
CORE-BEHRT, MOTOR and SurvivEHR have shown for structured records
\citep{li2020behrt,rasmy2021med,odgaard2024core,steinberg2023motor,gadd2026survivehr}
and Delphi-2M, life2vec and pan-cancer models across diseases and whole
populations
\citep{shmatko2025delphi,savcisens2024life2vec,placido2023pancreatic,Jee2024MSKCHORD,bjerregaard2025pancancer}.
In oncology, a mortality model shared across eleven cancer types improved when
it was then adapted to each type \citep{Oh2021twostage}, and studies pooling
lymphoid cancers have each examined a single outcome
\citep{parviz2026dacll,werling2026dlbcl}. But sharing does not help by default.
Multi-task studies of health records report uneven benefits and negative
transfer \citep{harutyunyan2019multitask,mcdermott2020mtl}, and joint learning
has helped rare phenotypes while hurting common ones \citep{ding2019phenotyping}.
Models that share experience across tasks
\citep{zhang2019metapred,liu2023taml,chandak2026everyquery} are also rarely
evaluated on outcomes that matter to health systems \citep{wornow2023shaky} or
on how few labels they need \citep{wornow2023ehrshot}. Less is known about
where shared signal comes from when many related tasks are learned together, how
much each source contributes against the tools in clinical use, and whether a
model can learn to take what helps and leave what does not.

Haematological cancers are a fitting place to study this. Lymphoid cancers
comprise dozens of entities \citep{alaggio2022who}, each uncommon
\citep{sant2010haemacare}: in the cohort studied here, 24 subtypes range from
almost 10,000 patients with diffuse large B-cell lymphoma (DLBCL) to fewer than
50 with amyloidosis. Even the most common subtype offers few events for most
questions. A complication with a one-year risk of 2\% gives fewer than 200
events among the patients with DLBCL treated before 2022, and in a subtype of
500 patients a risk of 5\% gives 25. Yet at the start of treatment a
haematologist weighs treatment failure alongside infection, cytopenia, organ
toxicity and unplanned admission. The tools for this follow the one-task
template. Prognostic indices such as the International Prognostic Index (IPI)
and its NCCN refinement, the CLL-IPI and the revised International Staging
System turn a handful of measurements at diagnosis into points for one disease
and one endpoint \citep{IPI1993,zhou2014nccnipi,cllipi2016,palumbo2015riss},
losing information that the same measurements carry when kept continuous
\citep{biccler2018simplicity}. They matter nonetheless: the Danish DLBCL
guideline sets the number of cycles and the intensity of treatment by the
age-adjusted IPI \citep{DLG2025DLBCL}, and indices decide entry to trials of new
first-line therapy \citep{Tilly2022POLARIX,Langerbeins2022CLL12}. Machine
learning on the wider record has improved on them, one problem at a time
\citep{Biccler2018,agius2020clltim,werling2026dlbcl}. Under the NCCN-IPI no
patient over 60 can be placed at low risk, whereas a model reading the record
placed 160 such patients with DLBCL at low risk of treatment failure, against
one by the index, with comparable outcomes \citep{werling2026dlbcl}. For many
complications and subtypes, no dedicated predictor exists. The DALY-CARE
resource was built partly on the expectation that modelling the lymphoid
cancers together would improve prediction for each of them
\citep{brieghel2025dalycare}, and Denmark makes that expectation testable: a personal identifier links hospital contacts, prescriptions,
laboratory results and deaths across the population
\citep{schmidt2014cpr,schmidt2019danish}, and national quality registries record
staging, treatment and response with near-complete coverage
\citep{Arboe2016lyfo,daCunhaBang2016,Gimsing2016}.

We expected sharing to help broadly, because every task is short of events, and
most where a prediction has least data of its own. Where a target has more
data, sharing can as easily dilute what is specific to it as add to it. Learning
what to borrow addresses this: a model that chooses its sources should keep the
gains of sharing where they exist and avoid its costs where they do not.

The broadest source of signal is the population itself. An encoder pretrained
on the records of every resident has learned from diseases, treatments and
care far beyond haematology, and it represents records in finer detail: codes
too rare to learn in a cohort, such as specific molecular findings, keep their
own token at national scale. If such an encoder can be reused for specialised
prediction, the investment in learning from population records could serve
many clinical questions. Continued pretraining on the target population,
domain adaptation, is the established way to close remaining gaps
\citep{gururangan2020dont,guo2024multicenter,lemmon2023pediatric}; here the
national population already contains the cohort, so adaptation adds no
patients and only refocuses the encoder.

In a nationwide cohort of 37,531 patients with 24 lymphoid cancer subtypes, we
predicted 69 outcomes at six horizons from the start of first-line treatment and
widened the sources of signal step by step (\cref{fig:paradigms}): from one
disease and one outcome, to all diseases, all outcomes or both, to a model that
learns what to borrow from each disease and outcome, and finally to an encoder
pretrained on the whole population. Each step was compared on the same patients
under one protocol, against gradient-boosted trees tuned for every prediction
and the prognostic indices recorded in the registers, in a single evaluation of
patients treated after every modelling decision had been fixed. We ask how much
each source of shared signal adds, where it helps most, and whether letting the
model choose its sources is what makes sharing work.

\clearpage
\begin{figure}[H]
  \centering
  \resizebox{\textwidth}{!}{\begin{tikzpicture}[opera,x=1cm,y=1cm]
  \draw[divider] (5.55,.05)--(5.55,-6.18);
  \draw[divider] (11.15,.05)--(11.15,-6.18);

  % Shared semantic bands make the three paradigms directly comparable.
  \begin{scope}[on background layer]
    \draw[draw=OperaBorder,densely dashed,line width=.45pt] (-.42,-2.70)--(16.55,-2.70);
    \draw[draw=OperaBorder,densely dashed,line width=.45pt] (-.42,-4.72)--(16.55,-4.72);
  \end{scope}
  \node[eyebrow,anchor=east,text=OperaInkMuted] at (-.54,-2.15) {INPUT};
  \node[eyebrow,anchor=east,text=OperaInkMuted] at (-.54,-3.68) {MODEL};
  \node[eyebrow,anchor=east,text=OperaInkMuted] at (-.54,-5.62) {OUTPUT};

  % A — clinical score
  \node[panelletter,anchor=north west] at (0,0) {A};
  \node[eyebrow,anchor=north west] at (.48,-.03) {CLINICAL SCORE};
  \node[heading,anchor=north west] at (.48,-.36) {A diagnostic snapshot};
  \node[labelsmall,anchor=north west,align=left,text width=4.65cm] at (.48,-.78)
    {Few variables, fixed combination, one disease};
  \foreach \name/\x/\lab in {age/.72/Age,stage/1.68/Stage,
      ldh/2.65/LDH,ps/3.62/PS,sites/4.58/Sites}
    \node[circle,draw=OperaSienna!55,fill=OperaSienna!7,minimum size=7.4mm,
      font=\sffamily\scriptsize] (\name) at (\x,-2.15) {\lab};
  \node[stagewarm,minimum width=2.15cm,minimum height=7mm] (score) at (2.65,-3.70) {risk score};
  \foreach \name in {age,stage,ldh,ps,sites}
    \draw[-{Latex[length=1.5mm]},draw=OperaSienna,line width=.65pt]
      (\name.south)--(score.north);
  \draw[flowwarm,shorten >=1pt] (score.south)--(2.65,-5.25);
  % Clinical scores define discrete strata, not a continuous probability scale.
  \node[labelsmall] at (1.38,-5.48) {low};
  \node[labelsmall] at (2.65,-5.48) {intermediate};
  \node[labelsmall] at (3.92,-5.48) {high};
  \fill[OperaSienna!16] (.75,-6.00) rectangle (2.02,-5.72);
  \fill[OperaSienna!43] (2.02,-6.00) rectangle (3.28,-5.72);
  \fill[OperaSienna!82] (3.28,-6.00) rectangle (4.55,-5.72);
  \draw[white,line width=.8pt] (2.02,-6.00)--(2.02,-5.72);
  \draw[white,line width=.8pt] (3.28,-6.00)--(3.28,-5.72);
  %\node[captionnote,text width=4.4cm] at (2.65,-6.72)
  %  {Interpretable and portable, but most prior history is compressed away.};

  % B — tabular model
  \node[panelletter,anchor=north west] at (5.86,0) {B};
  \node[eyebrow,anchor=north west] at (6.34,-.03) {TABULAR MACHINE LEARNING};
  \node[heading,anchor=north west] at (6.34,-.36) {History becomes features};
  \node[labelsmall,anchor=north west,align=left,text width=4.55cm] at (6.34,-.78)
    {Longitudinal history, engineered summaries, one task};
  \draw[OperaInkMuted,line width=.55pt] (6.35,-2.15)--(10.38,-2.15);
  \foreach \x/\col in {6.55/OperaPeriwinkle,7.22/OperaTeal,7.95/OperaTerracotta,
      8.63/OperaTeal,9.42/OperaOchre,10.08/OperaPeriwinkle}
    \fill[\col] (\x,-2.15) circle (2.1pt);
  % One history, partitioned into look-back intervals for aggregation.
  \foreach \x in {6.55,8.42,9.56,10.18}
    \draw[draw=OperaTeal!68,line width=.65pt] (\x,-2.03)--(\x,-2.27);
  \node[font=\sffamily\tiny,text=OperaTeal] at (7.48,-2.43) {365 d};
  \node[font=\sffamily\tiny,text=OperaTeal] at (8.99,-2.43) {180 d};
  \node[font=\sffamily\tiny,text=OperaTeal] at (9.87,-2.43) {30 d};
  \node[stage,draw=OperaTeal!65,fill=OperaTeal!7,minimum width=2.80cm,
    minimum height=7mm] (ml) at (8.36,-3.70)
    {ML algorithm\\{\scriptsize e.g. XGBoost}};
  \draw[flowteal,shorten >=1pt] (8.36,-2.42)--(ml.north);
  \node[labelsmall] (prob) at (8.36,-5.48) {$P(\mathrm{outcome}_n\ \mathrm{by}\ t)$};
  \draw[flowteal,shorten >=1pt] (ml.south)--(prob.north);
  \foreach \i in {0,...,23}{
    \pgfmathsetmacro{\xa}{6.80+0.13*\i}
    \pgfmathsetmacro{\xb}{6.80+0.13*(\i+1)}
    \pgfmathtruncatemacro{\tone}{6+3*\i}
    \fill[OperaTeal!\tone] (\xa,-6.00) rectangle (\xb,-5.72);
  }
  %\node[captionnote,text width=4.4cm] at (8.36,-6.72)
  %  {A separate predictor is typically fitted for each outcome and horizon.};

  % C — sequence model
  \node[panelletter,anchor=north west] at (11.46,0) {C};
  \node[eyebrow,anchor=north west] at (11.94,-.03) {LONGITUDINAL SEQUENCE MODEL};
  \node[heading,anchor=north west] at (11.94,-.36) {The timeline remains visible};
  \node[labelsmall,anchor=north west,align=left,text width=4.75cm] at (11.94,-.78)
    {Ordered events, contextual states, reusable learning};
  \draw[OperaIndigo,line width=.65pt] (11.95,-2.15)--(16.25,-2.15);
  \foreach \x/\col/\lab in {12.12/OperaPeriwinkle/D,12.82/OperaTeal/L,
      13.52/OperaTerracotta/M,14.22/OperaOchre/V,14.92/OperaTeal/L,15.62/OperaLavender/P}
    \node[circle,draw=\col,fill=white,minimum size=5.8mm,font=\sffamily\tiny] at (\x,-2.15) {\lab};
  \draw[flowblue] (14.10,-2.51)--(14.10,-2.80);
  \node[stageblue,minimum width=3.00cm,minimum height=7mm] (enc) at (14.10,-3.23)
    {transformer encoder};
  \node[softpanel,draw=OperaIndigo!42,fill=OperaIndigo!4,
    minimum width=2.25cm,minimum height=5.5mm,font=\sffamily\tiny]
    (z) at (14.10,-4.20) {shared representation $z$};
  \draw[flowblue,shorten >=1pt] (enc.south)--(z.north);
  \node[labelsmall] (c1) at (12.20,-5.48) {$P($tox. by $t)$};
  \node[labelsmall] (c2) at (13.70,-5.48) {$P($anemia by $t)$};
  \node[font=\sffamily\scriptsize,text=OperaInkMuted] at (14.80,-5.48) {$\cdots$};
  \node[labelsmall] (cn) at (16.02,-5.48) {$P($outcome$_n$ by $t)$};
  \draw[flowblue,shorten >=1pt] (z)--(c1.north);
  \draw[flowblue,shorten >=1pt] (z)--(c2.north);
  \draw[flowblue,shorten >=1pt] (z)--(cn.north);
  % Each task emits its own continuous probability, aligned with panel B.
  \begin{scope}[shift={(11.72,0)}]
    \foreach \i in {0,...,11}{
      \pgfmathsetmacro{\xa}{0.08*\i}
      \pgfmathsetmacro{\xb}{0.08*(\i+1)}
      \pgfmathtruncatemacro{\tone}{6+7*\i}
      \fill[OperaIndigo!\tone] (\xa,-6.00) rectangle (\xb,-5.72);
    }
  \end{scope}
  \begin{scope}[shift={(13.22,0)}]
    \foreach \i in {0,...,11}{
      \pgfmathsetmacro{\xa}{0.08*\i}
      \pgfmathsetmacro{\xb}{0.08*(\i+1)}
      \pgfmathtruncatemacro{\tone}{6+7*\i}
      \fill[OperaIndigo!\tone] (\xa,-6.00) rectangle (\xb,-5.72);
    }
  \end{scope}
  \begin{scope}[shift={(15.54,0)}]
    \foreach \i in {0,...,11}{
      \pgfmathsetmacro{\xa}{0.08*\i}
      \pgfmathsetmacro{\xb}{0.08*(\i+1)}
      \pgfmathtruncatemacro{\tone}{6+7*\i}
      \fill[OperaIndigo!\tone] (\xa,-6.00) rectangle (\xb,-5.72);
    }
  \end{scope}
  %\node[captionnote,text width=4.55cm] at (14.10,-6.72)
  %  {The encoder may be pretrained; one representation can support many outcomes and horizons.};
\end{tikzpicture}
}\par\vspace{3mm}
  \resizebox{\textwidth}{!}{\begin{tikzpicture}[opera,x=1cm,y=1cm]
  \draw[divider] (7.25,.05)--(7.25,-9.45);

  \node[panelletter,anchor=north west] at (0,0) {D};
  \node[eyebrow,anchor=north west] at (.48,-.03) {SOURCES OF PREDICTIVE INFORMATION};
  \node[heading,anchor=north west] at (.48,-.38) {More information for one prediction problem};
  \node[labelsmall,anchor=north west,align=left,text width=5.9cm] at (.48,-.82)
    {Moving right adds information to the learning problem.};

  % Concrete before-to-after examples are easier to read than abstract axes.
  \node[eyebrow,text=OperaIndigo,anchor=east] at (1.58,-1.90) {PATIENTS};
  \foreach \x/\y in {1.92/-1.76,2.12/-2.03,2.35/-1.80}
    \fill[OperaIndigo!55] (\x,\y) circle (1.8pt);
  \draw[flowblue] (2.62,-1.90)--(3.42,-1.90);
  \foreach \x/\y in {3.72/-1.70,3.93/-2.02,4.15/-1.78,4.38/-2.07,4.58/-1.69,
      4.82/-1.94,5.05/-1.74,5.28/-2.06,5.52/-1.84}
    \fill[OperaIndigo!65] (\x,\y) circle (1.8pt);
  \node[labelsmall,anchor=west] at (1.82,-2.27) {larger sample};

  \node[eyebrow,text=OperaTeal,anchor=east] at (1.58,-3.15) {VARIABLES};
  \draw[OperaTeal!55,line width=.55pt] (1.86,-3.41) rectangle (2.52,-2.89);
  \draw[OperaTeal!55,line width=.45pt] (2.18,-3.41)--(2.18,-2.89);
  \draw[flowteal] (2.72,-3.15)--(3.42,-3.15);
  \draw[OperaTeal!65,line width=.55pt] (3.70,-3.51) rectangle (5.48,-2.79);
  \foreach \x in {4.05,4.40,4.75,5.10}\draw[OperaTeal!45,line width=.4pt] (\x,-3.51)--(\x,-2.79);
  \node[labelsmall,anchor=west] at (1.82,-3.72) {richer phenotype};

  \node[eyebrow,text=OperaSienna,anchor=east] at (1.58,-4.45) {TIME};
  \fill[OperaSienna!65] (2.18,-4.45) circle (2.2pt);
  \draw[flowwarm] (2.72,-4.45)--(3.42,-4.45);
  \draw[OperaSienna!70,line width=.65pt] (3.70,-4.45)--(5.48,-4.45);
  \foreach \x in {3.82,4.18,4.55,4.94,5.34}\fill[OperaSienna!70] (\x,-4.45) circle (1.9pt);
  \node[labelsmall,anchor=west] at (1.82,-4.85) {longer history};

  \draw[draw=OperaBorder,densely dashed,line width=.5pt] (1.82,-5.16)--(5.62,-5.16);

  \node[eyebrow,text=OperaIndigo,anchor=east] at (1.58,-5.72) {DISEASES};
  \foreach \x/\y/\col in {1.98/-5.63/OperaIndigo,2.28/-5.81/OperaPeriwinkle}
    \fill[\col!70] (\x,\y) circle (2.3pt);
  \draw[flowblue] (2.72,-5.72)--(3.42,-5.72);
  \foreach \x/\y/\col in {3.72/-5.58/OperaIndigo,4.05/-5.82/OperaPeriwinkle,
      4.38/-5.62/OperaTeal,4.72/-5.84/OperaLavender,5.05/-5.57/OperaIndigo,
      5.38/-5.79/OperaTeal,5.66/-5.62/OperaPeriwinkle}
    \fill[\col!72] (\x,\y) circle (2.3pt);
  \node[labelsmall,anchor=west] at (1.82,-6.08) {broader cohort space};

  \node[eyebrow,text=OperaTeal,anchor=east] at (1.58,-6.82) {OUTCOMES};
  \foreach \x/\y in {2.02/-6.72,2.32/-6.92}
    \fill[OperaTeal!55] (\x-.08,\y-.08) rectangle (\x+.08,\y+.08);
  \draw[flowteal] (2.72,-6.82)--(3.42,-6.82);
  \foreach \x/\y/\tone in {3.72/-6.68/38,4.05/-6.92/48,4.38/-6.72/58,
      4.72/-6.94/68,5.05/-6.67/44,5.38/-6.88/76,5.66/-6.70/56}
    \fill[OperaTeal!\tone] (\x-.08,\y-.08) rectangle (\x+.08,\y+.08);
  \node[labelsmall,anchor=west] at (1.82,-7.18) {more learning tasks};
  \node[eyebrow,text=OperaGreen,anchor=east] at (1.58,-7.48) {MORE INFORMATION};
  \draw[-{Latex[length=2.2mm]},draw=OperaGreen,line width=1pt] (1.82,-7.48)--(5.48,-7.48);
  \node[labelsmall,text=OperaGreen] at (3.65,-7.76)
    {higher attainable performance};

  \node[panelletter,anchor=north west] at (7.60,0) {E};
  \node[eyebrow,anchor=north west] at (8.08,-.03) {LEARNING ACROSS COHORTS AND TASKS};
  \node[heading,anchor=north west] at (8.08,-.38) {Related diseases and outcomes};
  \node[labelsmall,anchor=north west,align=left,text width=7.6cm] at (8.08,-.82)
    {More disease cohorts and more outcome tasks add structured information around a low-support target.};

  % Separate the axis/title and matrix from the two-line panel introduction.
  \begin{scope}[yshift=-5mm]
  % Keep the axis title separate from the column labels and first row.
  \draw[-{Latex[length=1.8mm]},draw=OperaIndigo,line width=.75pt]
    (10.36,-1.55)--(16.72,-1.55);
  \node[eyebrow,text=OperaIndigo] at (13.54,-1.30) {MORE OUTCOME TASKS};

  \foreach \r/\disease in {0/DLBCL,1/Multiple myeloma,2/CLL,3/Hodgkin lymphoma,4/PTCL}
    \node[labelsmall,anchor=east] at (10.20,-2.72-.84*\r) {\disease};
  \draw[-{Latex[length=1.8mm]},draw=OperaTeal,line width=.75pt]
    (7.76,-2.40)--(7.76,-7.62);
  \node[eyebrow,text=OperaTeal,rotate=90] at (7.47,-5.01) {MORE DISEASE COHORTS};
  \foreach \col/\outcome in {0/Anaemia,1/Infection,2/Renal,3/Fracture,4/Survival}
    \node[font=\sffamily\fontsize{6.2}{7}\selectfont,text=OperaInkSecondary,
      align=center,text width=1.05cm,anchor=south] at (10.82+1.12*\col,-2.27) {\outcome};
  \node[font=\sffamily\fontsize{6.2}{7}\selectfont,text=OperaInkSecondary,
    align=center,text width=1.05cm,anchor=south] at (16.72,-2.27) {Outcome$_n$};
  \foreach \r in {0,...,4}{\foreach \col in {0,...,4}{
    \pgfmathtruncatemacro{\tone}{7+3*\r+2*\col}
    \node[rounded corners=.35mm,draw=OperaBorder,fill=OperaIndigo!\tone,
      minimum width=8.8mm,minimum height=5.7mm] (m\r\col)
      at (10.82+1.12*\col,-2.72-.84*\r) {};
  }}
  \foreach \r in {0,...,4}{
    \node[font=\sffamily\scriptsize,text=OperaInkMuted]
      at (16.01,-2.72-.84*\r) {$\cdots$};
  }
  \foreach \r in {0,...,4}{
    \pgfmathtruncatemacro{\tone}{16+3*\r}
    \node[rounded corners=.35mm,draw=OperaBorder,fill=OperaIndigo!\tone,
      minimum width=8.8mm,minimum height=5.7mm] (mn\r) at (16.72,-2.72-.84*\r) {};
  }
  \foreach \x in {10.82,11.94,13.06,14.18,15.30,16.72}{
    \node[font=\sffamily\scriptsize,text=OperaInkMuted] at (\x,-6.60) {$\vdots$};
  }
  \node[font=\sffamily\scriptsize,text=OperaInkMuted] at (16.01,-6.60) {$\ddots$};
  \node[labelsmall,anchor=east] at (10.20,-7.29) {Disease$_m$};
  \foreach \col in {0,...,4}{
    \pgfmathtruncatemacro{\tone}{22+2*\col}
    \node[rounded corners=.35mm,draw=OperaBorder,fill=OperaIndigo!\tone,
      minimum width=8.8mm,minimum height=5.7mm] (mm\col)
      at (10.82+1.12*\col,-7.29) {};
  }
  \node[rounded corners=.35mm,draw=OperaBorder,fill=OperaIndigo!34,
    minimum width=8.8mm,minimum height=5.7mm] (mmn) at (16.72,-7.29) {};
  \node[rounded corners=.35mm,draw=OperaSienna,line width=1.1pt,fill=OperaSienna!8,
    minimum width=8.8mm,minimum height=5.7mm] (target) at (11.94,-6.08) {};
  % Every training cell is connected; most links recede behind the emphasized
  % same-outcome and same-disease borrowing paths.
  \begin{scope}[on background layer]
    \foreach \r in {0,...,4}{\foreach \col in {0,...,4}{
      \draw[draw=OperaTeal,opacity=.55,line width=.6pt] (m\r\col)--(target);
    }}
    \foreach \r in {0,...,4}{
      \draw[draw=OperaTeal,opacity=.55,line width=.6pt] (mn\r)--(target);
    }
    \foreach \col in {0,...,4}{
      \draw[draw=OperaTeal,opacity=.55,line width=.6pt] (mm\col)--(target);
    }
    \draw[draw=OperaTeal,opacity=.55,line width=.6pt] (mmn)--(target);
    \foreach \r in {0,...,3}{
      \draw[draw=OperaTeal!85!black,opacity=1,line width=1.6pt] (m\r1)--(target);
    }
    \foreach \col in {0,2,3,4}{
      \draw[draw=OperaTeal!85!black,opacity=1,line width=1.6pt] (m4\col)--(target);
    }
    % The last disease (same outcome) and last outcome (same disease) borrow too.
    \draw[draw=OperaTeal!85!black,opacity=1,line width=1.6pt] (mm1)--(target);
    \draw[draw=OperaTeal!85!black,opacity=1,line width=1.6pt] (mn4)--(target);
  \end{scope}
  \end{scope}
\end{tikzpicture}
}
  \caption{\textbf{Approaches to risk prediction and the information each
  can use.} Top, three approaches for haematological risk prediction.
  \textbf{A}, clinical risk scores combine a small diagnostic snapshot with a
  fixed rule into broad risk strata. \textbf{B}, classical machine learning
  uses a wider history, aggregated over time into a task-specific feature
  table. \textbf{C}, a longitudinal sequence model reads the ordered clinical
  events and can share one patient representation across outcomes; its
  encoder may be pretrained on the cohort or on the whole population. Bottom,
  sources of additional signal. \textbf{D}, conventional scaling adds patients,
  variables or observations to one isolated prediction problem. \textbf{E},
  shared learning also draws on related diseases (rows) and outcomes
  (columns).}
  \label{fig:paradigms}
\end{figure}

\clearpage
\operasection{Methods}

\operasubsection{Population, prediction time and outcomes}
We included adults in DALY-CARE who started first-line treatment for a lymphoid
or plasma-cell malignancy, using linked Danish population registers and clinical
quality registries \citep{brieghel2025dalycare}. The models read pre-treatment
diagnoses, procedures, pathology, medication, laboratory results and hospital
contacts, with age and sex (\cref{fig:model-inputs}); quality-registry fields
were used for cohort definitions, outcomes and clinical benchmarks, not as
inputs. The study was approved by the Danish Health Data Authority (P-2020-561)
and the National Ethics Committee (1804410).

Every prediction was made at the start of first-line treatment, from records
dated strictly before it, to support prognosis, monitoring and supportive-care
planning. We predicted 69 outcomes at 30, 60, 90, 180, 365 and 730 days,
spanning disease course, infection, laboratory abnormalities, cardiovascular
events and acute care, with death as a competing event for every non-fatal
outcome. Laboratory outcomes used CTCAE-aligned thresholds
\citep{nci2025ctcae}, and a patient contributed to an outcome only where it
could be ascertained, for example with laboratory monitoring. The subtype was
assigned from diagnoses recorded up to 14 days after treatment started, or 90
days for diagnoses coded late, as the diagnostic work-up is usually complete by
then; a sensitivity analysis used only diagnoses recorded before treatment.
Treatment lines, subtype assignment, code lists and ascertainment are detailed
in the Supplementary Methods.

\operasubsection{Temporal design}
Patients treated through 2021 formed the training set, those treated in 2022
the tuning set, and those treated from 2023 the held-out test set, which was
evaluated once, after every modelling decision had been fixed. Pretraining used
events before 1 January 2022. The analysis plan was locked before the test set
was opened \placeholder{date; supplementary analysis plan}.

\operasubsection{Learning across tasks}
All sequence models share one architecture: a causal transformer, building on
CORE-BEHRT \citep{odgaard2024core}, pretrained to predict each next event and
fine-tuned with a head that gives the cumulative incidence of every outcome at
every horizon, with death competing \citep{gensheimer2019nnetsurvival,lee2018deephit}.
Only the sources of signal differ (\cref{fig:supervision-grid}). A specialist
learns one outcome in one disease. A disease-pooled model learns one outcome
across all diseases, an outcome-pooled model all outcomes within one disease,
and the joint model every outcome in every disease. These four scopes separate
the signal gained from other diseases from that gained from other outcomes.

The joint model shares everything with every prediction. The routed model
instead learns what to borrow: each disease and each outcome has a small adapter
trained only on its own labels, and for every patient and target outcome a
router weighs these adapters, so that a prediction can draw on the diseases and
outcomes that help it and ignore the rest (\cref{fig:source-router})
\citep{ma2018mmoe,houlsby2019adapters,hu2022lora}. Three controls separate
learning what to borrow from simpler explanations: a router without adapters
tests added capacity, a joint model told the patient's disease tests knowing the
diagnosis, and models restricted to the patient's own disease or outcome adapter
test task-specific capacity without borrowing.

The broadest source is the population. Joint and routed models were also
initialised from the PHAIR encoder \placeholder{PHAIR reference}, pretrained
on the registers of the whole Danish population by its developers, either directly or after continued pretraining on the
cohort's records. Because the national population contains the cohort, national
against cohort pretraining measures what the rest of the population adds, and
adaptation what refocusing adds \citep{gururangan2020dont}. To test whether
population pretraining reduces the labels needed, the joint model was retrained
from both encoders on nested 10\%, 25\% and 50\% subsets of the training
patients, with the full tuning and test sets; against specialists trained on
all patients, the same subsets show how much labelled data sharing saves. Random initialisation and a frozen
encoder serve as references for pretraining and fine-tuning. Architecture,
training and the national data are described in the Supplementary Methods.

The configuration that performed best on the tuning set, across encoder origin
and routing, is called \opera{} (\modelname) and is the model compared with
clinical practice \placeholder{selection metric}. The contrasts between scopes,
routing and encoders use their fixed arms.

\operasubsection{Comparators}
Gradient-boosted trees remain hard to beat on tabular data
\citep{grinsztajn2022tabular}, so XGBoost \citep{chen2016xgboost} and ridge
logistic regression were tuned for
every disease, outcome and horizon on summaries of the same register histories,
built within each disease with the pipeline of our earlier DLBCL study
\citep{werling2026dlbcl,bernstorff2023tsflattener}. XGBoost pooled across
diseases, on features built across all patients, tests whether pooling also
helps tabular prediction, and XGBoost given the quality-registry variables
shows what curated clinical fields add. Clinical indices were compared with
\opera{} and the tabular models on the patients with the index recorded, on the
endpoint each index was developed for. The IPI in DLBCL is the main example,
with the NCCN-IPI and every other index recorded for at least 50 patients of a
subtype \citep{IPI1993,zhou2014nccnipi}; model risk groups were fixed on tuning
data to match the index's low-risk observed risk and high-risk specificity.

\operasubsection{Evaluation}
One-year time-dependent AUC was the primary metric, with the index of
prediction accuracy (IPA), Brier score and calibration for probability accuracy
\citep{blanche2013competingauc,gerds2006brier,schoop2011competing,kattan2018ipa}.
Every contrast was paired on the same patients, with intervals from 1,000
patient resamples within disease shared across models, outcomes and horizons.
Each cell's paired effect is reported alongside hierarchical summaries that
estimate the typical effect, its variation between diseases and outcomes, and
its relation to the number of training events. For the DLBCL example, decision
curves compare the net benefit of acting on \opera{}, the IPI and XGBoost
when patients above a predicted two-year risk of death receive closer monitoring
and supportive care, across thresholds from 5\% to 40\%, fixed before the test
set was opened \citep{vickers2006dca} \placeholder{confirm action and range}. Training time, compute and storage are reported for one
\opera{} model and for the separately trained models it replaces.
Integrated gradients \citep{sundararajan2017ig} describe what \opera{} relies
on, by source, time and clinical finding and for individual patients, checked by
event removal and weight randomisation \citep{adebayo2018sanity}. Reporting
follows TRIPOD+AI \citep{collins2024tripodai}.

% Result macros (T2, T9, ipi_*, ig_*) must be computed for the selected OPERA
% configuration; the scope contrasts (P*, R*, E*) use the fixed arms.
\FloatBarrier
\operasection{Results}

\operasubsection{Who was studied, and what was predicted?}
The pretraining population comprised all 78,668 adults in DALY-CARE with a
lymphoid cancer diagnosis; their records before 2022 pretrained the cohort
encoder. Of these, 38,635 started first-line treatment, and the 37,531 with an
assigned lymphoid disease who were alive and followed at treatment start form
the final cohort, on which every prediction model was trained,
tuned and tested (\cref{fig:cohort-flow}): 29,105 in the training period (to
2021), 1,980 in the tuning year (2022) and 6,446 in the held-out test set (from
2023). The median age at treatment was 69 years (IQR 60--77), 41.7\% were
women, and the largest diseases were diffuse large B-cell lymphoma (26.5\%),
multiple myeloma (23.0\%), chronic lymphocytic leukaemia (10.9\%) and
follicular lymphoma (7.8\%); patients in the test period were older (median 72
years; \cref{tab:characteristics}). Of the 24 diseases and 69 outcomes,
\result{cells_evaluable}{N} cells had enough test events to be evaluated at one year. Event rates ranged from \placeholder{range}, and \result{cells_low_support}{N} cells had fewer than 20 events to train on
(\cref{fig:supp-coverage,fig:supp-support}). The tuning set selected
\placeholder{configuration} as \opera{}.

% Only the patient-count row below has been updated for the released supervised population.
% All other numeric rows remain historical layout values and MUST be re-exported before use.
\begin{table*}[t]
\centering
\caption{Baseline characteristics of the supervised first-line haematology analysis cohort. \placeholder{All characteristics except patient counts require re-export from the current release.}}
\label{tab:characteristics}
\footnotesize
\setlength{\tabcolsep}{4pt}
\begin{tabularx}{\textwidth}{@{}Xrrrr@{}}
\toprule
Characteristic & Overall & Training & Tuning & Held-out \\
\midrule
Patients, $n$ & 37,531 & 29,105 & 1,980 & 6,446 \\
Age, years, median (IQR) & 69.3 (59.7--77.1) & 68.7 (59.2--76.5) & 71.1 (61.1--78.4) & 71.7 (61.7--78.8) \\
Age $\geq$70 years, $n$ (\%) & 18,011 (48.1\%) & 13,431 (46.1\%) & 1,050 (53.4\%) & 3,530 (55.2\%) \\
Female, $n$ (\%) & 15,647 (41.8\%) & 12,199 (41.9\%) & 802 (40.8\%) & 2,646 (41.4\%) \\
Male, $n$ (\%) & 21,820 (58.2\%) & 16,909 (58.1\%) & 1,165 (59.2\%) & 3,746 (58.6\%) \\
Year of first-line treatment, median (IQR) & 2016 (2010--2021) & 2014 (2009--2018) & 2022 (2022--2022) & 2024 (2023--2025) \\
Observed follow-up, years, median (IQR) & 4.1 (1.4--8.6) & 5.8 (2.1--10.1) & 3.5 (3.2--3.8) & 1.3 (0.6--2.2) \\
\addlinespace
\multicolumn{5}{@{}l}{\sffamily\bfseries Disease groups, $n$ (\%)}\\
Amyloidosis & 47 (0.1\%) & 22 (0.1\%) & 7 (0.4\%) & 18 (0.3\%) \\
\textbf{Burkitt/lymphoblastic lymphoma} & 1,364 (3.6\%) & 1,075 (3.7\%) & 87 (4.4\%) & 202 (3.2\%) \\
\hspace*{1em}Burkitt lymphoma (BL) & 295 (0.8\%) & 238 (0.8\%) & 15 (0.8\%) & 42 (0.7\%) \\
\hspace*{1em}Lymphoblastic lymphoma (LBL) & 1,069 (2.9\%) & 837 (2.9\%) & 72 (3.7\%) & 160 (2.5\%) \\
\textbf{CLL/SLL} & 5,152 (13.8\%) & 4,092 (14.1\%) & 244 (12.4\%) & 816 (12.8\%) \\
\hspace*{1em}Chronic lymphocytic leukaemia (CLL) & 4,062 (10.8\%) & 3,220 (11.1\%) & 187 (9.5\%) & 655 (10.2\%) \\
\hspace*{1em}Small lymphocytic lymphoma (SLL) & 1,090 (2.9\%) & 872 (3.0\%) & 57 (2.9\%) & 161 (2.5\%) \\
\textbf{DLBCL and related aggressive B-cell lymphomas} & 10,408 (27.8\%) & 8,369 (28.8\%) & 505 (25.7\%) & 1,534 (24.0\%) \\
\hspace*{1em}Diffuse large B-cell lymphoma (DLBCL) & 9,951 (26.6\%) & 8,008 (27.5\%) & 481 (24.5\%) & 1,462 (22.9\%) \\
\hspace*{1em}Richter transformation (RT) & 192 (0.5\%) & 151 (0.5\%) & 14 (0.7\%) & 27 (0.4\%) \\
\hspace*{1em}Transformed follicular lymphoma & 265 (0.7\%) & 210 (0.7\%) & 10 (0.5\%) & 45 (0.7\%) \\
Hairy-cell leukaemia & 418 (1.1\%) & 293 (1.0\%) & 29 (1.5\%) & 96 (1.5\%) \\
Hodgkin lymphoma & 2,644 (7.1\%) & 2,080 (7.1\%) & 141 (7.2\%) & 423 (6.6\%) \\
\textbf{Indolent B-cell non-Hodgkin lymphoma} & 5,399 (14.4\%) & 3,895 (13.4\%) & 302 (15.4\%) & 1,202 (18.8\%) \\
\hspace*{1em}Follicular lymphoma (FL) & 2,932 (7.8\%) & 2,248 (7.7\%) & 150 (7.6\%) & 534 (8.4\%) \\
\hspace*{1em}B-cell lymphoma, unspecified (BCL) & 227 (0.6\%) & 87 (0.3\%) & 17 (0.9\%) & 123 (1.9\%) \\
\hspace*{1em}Lymphoplasmacytic lymphoma (LPL) & 1,570 (4.2\%) & 1,165 (4.0\%) & 100 (5.1\%) & 305 (4.8\%) \\
\hspace*{1em}Extranodal marginal-zone lymphoma (EMZL) & 500 (1.3\%) & 318 (1.1\%) & 33 (1.7\%) & 149 (2.3\%) \\
\hspace*{1em}Nodal marginal-zone lymphoma (NMZL) & 27 (0.1\%) & 5 (0.0\%) & $<$5 & 22 (0.3\%) \\
\hspace*{1em}Splenic marginal-zone lymphoma (SMZL) & 143 (0.4\%) & 72 (0.2\%) & $<$5 & 69 (1.1\%) \\
Mantle-cell lymphoma & 1,466 (3.9\%) & 1,116 (3.8\%) & 87 (4.4\%) & 263 (4.1\%) \\
\textbf{Myeloma and related plasma-cell disorders} & 8,953 (23.9\%) & 6,904 (23.7\%) & 491 (25.0\%) & 1,558 (24.4\%) \\
\hspace*{1em}Multiple myeloma (MM) & 8,629 (23.0\%) & 6,673 (22.9\%) & 464 (23.6\%) & 1,492 (23.3\%) \\
\hspace*{1em}Plasma-cell leukaemia (PCL) & 201 (0.5\%) & 146 (0.5\%) & 18 (0.9\%) & 37 (0.6\%) \\
\hspace*{1em}Solitary myeloma/plasmacytoma & 123 (0.3\%) & 85 (0.3\%) & 9 (0.5\%) & 29 (0.5\%) \\
\textbf{T-cell non-Hodgkin lymphoma} & 1,616 (4.3\%) & 1,262 (4.3\%) & 74 (3.8\%) & 280 (4.4\%) \\
\hspace*{1em}Angioimmunoblastic T-cell lymphoma (AITL) & 483 (1.3\%) & 351 (1.2\%) & 26 (1.3\%) & 106 (1.7\%) \\
\hspace*{1em}Anaplastic large-cell lymphoma (ALCL) & 410 (1.1\%) & 339 (1.2\%) & 12 (0.6\%) & 59 (0.9\%) \\
\hspace*{1em}Peripheral T-cell lymphoma (PTCL) & 522 (1.4\%) & 444 (1.5\%) & 17 (0.9\%) & 61 (1.0\%) \\
\hspace*{1em}T-cell lymphoma, unspecified (TCL) & 201 (0.5\%) & 128 (0.4\%) & 19 (1.0\%) & 54 (0.8\%) \\
\addlinespace
\multicolumn{5}{@{}l}{\sffamily\bfseries Recorded performance status, $n$ (\%)}\\
Performance status 0 & 14,662 (53.8\%) & 11,842 (53.2\%) & 858 (54.3\%) & 1,962 (57.9\%) \\
Performance status 1 & 8,262 (30.3\%) & 6,838 (30.7\%) & 434 (27.5\%) & 990 (29.2\%) \\
Performance status 2 & 2,557 (9.4\%) & 2,077 (9.3\%) & 187 (11.8\%) & 293 (8.6\%) \\
Performance status 3 & 1,273 (4.7\%) & 1,063 (4.8\%) & 88 (5.6\%) & 122 (3.6\%) \\
Performance status 4 & 481 (1.8\%) & 444 (2.0\%) & 13 (0.8\%) & 24 (0.7\%) \\
Missing performance status$^{a}$ & 10,232 (27.3\%) & 6,844 (23.5\%) & 387 (19.7\%) & 3,001 (46.9\%) \\
\addlinespace
\multicolumn{5}{@{}l}{\sffamily\bfseries Disease-specific prognostic risk category, $n$ (\%)}\\
Low & 3,989 (17.6\%) & 3,231 (17.4\%) & 243 (17.6\%) & 515 (18.9\%) \\
Intermediate & 10,513 (46.4\%) & 8,608 (46.4\%) & 639 (46.3\%) & 1,266 (46.5\%) \\
High / very high$^{b}$ & 8,161 (36.0\%) & 6,721 (36.2\%) & 497 (36.0\%) & 943 (34.6\%) \\
Missing risk category$^{a}$ & 14,804 (39.5\%) & 10,548 (36.2\%) & 588 (29.9\%) & 3,668 (57.4\%) \\
\addlinespace
\multicolumn{5}{@{}l}{\sffamily\bfseries Pre-treatment record summaries, median (IQR)$^{c}$}\\
Recorded clinical-history span, years & \placeholder{Re-export} & \placeholder{Re-export} & \placeholder{Re-export} & \placeholder{Re-export} \\
Clinical-coded records & \placeholder{Re-export} & \placeholder{Re-export} & \placeholder{Re-export} & \placeholder{Re-export} \\
Distinct diagnosis codes & 21.0 (13.0--32.0) & 20.0 (12.0--30.0) & 24.0 (15.0--37.0) & 25.0 (15.0--38.0) \\
Distinct medication codes & 21.0 (12.0--34.0) & 18.0 (10.0--30.0) & 32.0 (21.0--46.0) & 34.0 (22.0--49.0) \\
Laboratory-coded records & 197.0 (0.0--482.0) & 118.0 (0.0--346.0) & 483.0 (279.5--803.5) & 531.0 (302.0--887.0) \\
\bottomrule
\end{tabularx}
\vspace{1.5mm}
\begin{minipage}{\textwidth}\scriptsize
One record per patient in the historical 37,467-person review export; these values are not the current 37,557-person first-line source or its supervised analysis subset. Disease groups are from \texttt{cohort\_grouped}, not the separate, incomplete \texttt{Disease} field. Indented rows show fine cohorts within their grouped parent; one-to-one groups are shown once. Disease percentages use the full cohort/split, not the parent group. A dash denotes an absent row in the source export, not a confirmed zero. Detailed treatment categories and scoring systems are reported in the supplementary tables.
$^{a}$Performance status and disease-specific prognostic risk categories are sourced from quality registries, whose disease and calendar-period coverage is incomplete. Category percentages use patients with nonmissing values; missingness percentages use the full cohort/split. The later-period decline in availability must be re-estimated from the current release before interpretation.
$^{b}$High and very high are combined only for this descriptive overview. The very-high category occurs in CLL-IPI (114 patients overall); risk categories from different scoring systems are not equivalent clinical scales. Original categories and scoring-system identities are retained in the supplement.
$^{c}$MEDS-derived values are provisional pending re-export after excluding demographic tokens from clinical-history summaries. Laboratory entries count coded records, not unique tests or verified source measurements. Follow-up is observed duration, not a reverse-Kaplan--Meier estimate.
IQR, interquartile range. Primary-only small-cell masking in the source exports is not disclosure clearance; the complete publication bundle requires institutional review.
\end{minipage}
\end{table*}


\FloatBarrier
\operasubsection{Which risks can be anticipated at the start of treatment?}
\placeholder{At one year, \opera{} discriminated outcomes in the X and Y
domains well (AUC A--B) and those in Z poorly (AUC C--D)} (\cref{fig:atlas}).
\placeholder{Change in AUC from 30 days to two years, and variation across
diseases} (\cref{fig:supp-atlas-matrix}).

\begin{figure}[H]
  \resultfigure{fig_atlas}{Discrimination across outcomes}
    {Test-period AUC of OPERA for every outcome at six horizons,
    with clinical-domain summaries.}
  \caption{\textbf{Discrimination of \opera{} across outcomes and
  horizons.} \textbf{A}, Test-period AUC for each of the 69
  outcomes at each horizon (colour), against its training events.
  \textbf{B}, The same results summarised by clinical domain. Outcomes or
  horizons without enough test events are not shown, and every outcome in every disease is in \cref{fig:supp-atlas-matrix}.}
  \label{fig:atlas}
\end{figure}

\FloatBarrier
\operasubsection{Does one model outperform the methods in use?}
\opera{} had the higher one-year AUC in \result{T2_better}{X of Y} disease--outcome cells against per-task XGBoost (median difference \result{T2_median}{Z}), in \result{T9_better}{X of Y} against XGBoost given the quality-registry variables, and in \result{P5_better}{X of Y} against the specialist trained on each cell alone (\cref{fig:headline}A--B, \cref{tab:panel-effects}). Pooled over cells, its difference from per-task XGBoost was \result{T2_auc}{X} in AUC and \result{T2_ipa}{Y} in IPA. \placeholder{Its median gain
over XGBoost was positive / not at every horizon from 30 days to two years}
(\cref{fig:headline}D) \placeholder{and did / did not differ with the number of
events recorded before treatment} (\cref{fig:headline}E). \placeholder{At the
specificity of each clinical index's high-risk category, the model's high-risk
group contained X to Y percentage points more of the patients with the event}
(\cref{fig:headline}C).

\begin{table}[H]
\centering\small
\caption{\textbf{Panel-level comparisons at one year.} Paired differences on the held-out test set pooled with equal weight per cell, with 95\% paired bootstrap
intervals. Cells: disease--outcome cells estimable for both models. Rare subtypes: subtypes with fewer than 500 training patients. Cells better: share of cells with a higher AUC point estimate, irrespective of uncertainty. Bold: the interval lies above zero, favouring the first model. Red: the interval lies below zero, favouring the second. The hierarchical estimates, which
weight diseases and endpoints equally, and the prespecified labels are in
Supplementary Table~\placeholder{S}.}
\label{tab:panel-effects}
\resizebox{\textwidth}{!}{%
\begin{tabular}{@{}lrrrrrl@{}}
\toprule
Comparison & Cells & $\Delta$AUC (95\% CI) & $\Delta$AUC, rare subtypes & $\Delta$IPA (95\% CI) & Cells better & Hierarchical label \\
\midrule
\IfFileExists{\resultnumbersdir panel_effects_rows.tex}{\input{\resultnumbersdir panel_effects_rows.tex}}{%
Joint vs specialist & \placeholder{N} & \placeholder{x} & \placeholder{x} & \placeholder{x} & \placeholder{x\%} & \placeholder{} \\
Disease-pooled vs specialist & \placeholder{N} & \placeholder{x} & \placeholder{x} & \placeholder{x} & \placeholder{x\%} & \placeholder{} \\
Outcome-pooled vs specialist & \placeholder{N} & \placeholder{x} & \placeholder{x} & \placeholder{x} & \placeholder{x\%} & \placeholder{} \\
Routed vs joint & \placeholder{N} & \placeholder{x} & \placeholder{x} & \placeholder{x} & \placeholder{x\%} & \placeholder{} \\
National vs cohort encoder (joint) & \placeholder{N} & \placeholder{x} & \placeholder{x} & \placeholder{x} & \placeholder{x\%} & \placeholder{} \\
Adapted national vs cohort encoder (joint) & \placeholder{N} & \placeholder{x} & \placeholder{x} & \placeholder{x} & \placeholder{x\%} & \placeholder{} \\
OPERA vs per-task XGBoost & \placeholder{N} & \placeholder{x} & \placeholder{x} & \placeholder{x} & \placeholder{x\%} & \placeholder{} \\
OPERA vs XGBoost pooled across diseases & \placeholder{N} & \placeholder{x} & \placeholder{x} & \placeholder{x} & \placeholder{x\%} & \placeholder{} \\
OPERA vs XGBoost with registry variables & \placeholder{N} & \placeholder{x} & \placeholder{x} & \placeholder{x} & \placeholder{x\%} & \placeholder{} \\
\bottomrule}
\end{tabular}}
\end{table}

\begin{figure}[H]
  \resultfigure{fig_headline}{OPERA versus reference methods}
    {Gains over each reference per cell, against XGBoost per cell, against the
    clinical indices, across horizons and by recorded history.}
  \caption{\textbf{Performance of \opera{} relative to reference
  methods.} \textbf{A}, One-year AUC gain over each
  reference, one point per disease--outcome cell with enough test events. Diamonds show medians, bars interquartile ranges, and the figures on the right the share of cells in which \opera{} was better. \textbf{B}, One-year AUC of \opera{}
  against per-task XGBoost in each cell, with point size reflecting training events.
  \textbf{C}, Share of patients with the event placed in the high-risk group by
  each clinical index and by \opera{} at the index's specificity (cutoffs
  fixed on tuning data). \textbf{D}, Median AUC gain over per-task XGBoost of
  each neural model at each horizon. \textbf{E}, Improvement in one-year Brier
  score over the specialist and over XGBoost by the number
  of events recorded before treatment ($\times$100, 95\% CI). Within each contrast, models are evaluated on the same test patients.}
  \label{fig:headline}
\end{figure}

Against XGBoost tuned for each disease, outcome and horizon on the same patients, \opera{}'s one-year AUC differed by \result{T2_auc}{X} and its IPA by \result{T2_ipa}{Y} (\cref{fig:tabular}). The specialists differed from XGBoost by \result{T3_auc}{X} in AUC, and XGBoost from logistic regression by \result{T4_auc}{Y}. XGBoost pooled across diseases differed from per-disease XGBoost by \result{T7_auc}{Y}, and \opera{} from pooled XGBoost by \result{T6_auc}{X}, which separates the contribution of pooling from that of the sequence model. \placeholder{Clinical domains and diseases in which XGBoost remained better.}
At a specificity of 90\% fixed on tuning data, \opera{}'s median sensitivity was \result{dp90_sens_joint}{X} against \result{dp90_sens_xgb}{Y} for XGBoost, and its median MCC \result{dp90_mcc_joint}{X} against \result{dp90_mcc_xgb}{Y}. \placeholder{Whether the ordering held at the other decision points.} Adding the quality-registry variables changed XGBoost's AUC by \result{T10_auc}{X}, and \opera{}, without them, differed from this reference by \result{T9_auc}{Y}.
\placeholder{Calibration slopes and calibration-in-the-large by model and
clinical domain} (\cref{fig:tabular}B, \cref{fig:calibration}).

\begin{figure}[H]
  \resultfigure{fig_benchmarks}{Comparison with tabular models}
    {A: AUC gains over XGBoost by clinical domain. B: calibration slope.
    C: gain against training events. D--E: reliability.}
  \caption{\textbf{Comparison with tabular models: discrimination and
  calibration.} \textbf{A}, One-year AUC gain over per-task XGBoost of the
  joint, specialist and routed models, and of \opera{} over XGBoost
  pooled across diseases and XGBoost given quality-registry variables, by
  clinical domain (median and interquartile range across cells).
  \textbf{B}, Calibration slope of the joint, routed, specialist and XGBoost
  models by domain (ideal 1). \textbf{C}, Gain of \opera{} over XGBoost
  against the target's training events, with the support curve. \textbf{D--E},
  Reliability of \opera{} and XGBoost for two example outcomes. Observed risks are censoring-weighted and point size reflects patients per
  bin. No model was recalibrated on test data. The full calibration analysis is in \cref{fig:calibration}.}
  \label{fig:tabular}
\end{figure}

One \opera{} model replaced \result{narrow_fits}{N} separately trained
specialists and \placeholder{N} tabular models, with \placeholder{training time,
GPU-hours and storage for one model against the portfolio it replaces}. Against
choosing, for each cell, whichever of the specialist, disease-pooled and
outcome-pooled models was best on the tuning set, the joint model's one-year
AUC differed by \result{J_vs_best_tuning_auc}{X} over
\result{J_vs_best_tuning_cells}{N} cells (\cref{fig:strategy}).

\FloatBarrier
\operasubsection{Does it improve on the clinical indices?}
\placeholder{N of N eligible patients with DLBCL had a recorded IPI, and how they differed from those without. Across subtypes, the share of patients with their subtype's index recorded or computable ranged from X\% to Y\% (Supplementary Table~S), whereas \opera{} gave an estimate for every patient.} Per 1,000 patients, using a threshold matched on tuning data to the IPI
high-risk category's specificity, \opera{} identified \result{ipi_extra_per1000}{N} more of the patients who died within two years, and placed \result{ipi_low_extra_per1000}{M} more at low risk, with an observed risk of \result{ipi_low_risk_model}{X\%} against \result{ipi_low_risk_index}{Y\%} in the IPI low-risk group (\cref{fig:risk-groups}). Its highest-risk 5\% had an observed two-year risk of \result{top5_risk_model}{X\%}, against \result{ipi_high_risk_index}{Y\%} in the IPI high-risk group. \placeholder{At that operating point, recall and
precision were X and Y for the model and A and B for the IPI, and average precision across thresholds was X versus Y} (\cref{fig:clinical-score}).
\placeholder{Outcomes of the patients the two systems classified differently,
including older patients whom the model placed at low risk}
(\cref{fig:risk-groups}D, \cref{fig:clinical-score-age}). \placeholder{Across thresholds of 5\% to 40\%, acting on \opera{}
gave a net benefit of X more true positives per 1,000 than the IPI and Y more
than XGBoost} (\cref{fig:decision-curve}).
\placeholder{Against the NCCN-IPI, result; XGBoost and logistic regression on the
same patients.} Median follow-up, by the reverse Kaplan--Meier method, was \result{followup_ipi}{X years (IQR)}. \placeholder{Across the other recorded indices, each in its own population and on its own endpoint, the model groups separated outcomes more / less clearly than the index in N of M} (\cref{fig:score-overview}).

\begin{figure}[H]
  \resultfigure{fig_risk_groups}{Risk groups versus the IPI}
    {A--B: overall survival by model risk group and by IPI group. C:
    two-year overall survival per group. D: cross-classification.}
  \caption{\textbf{Risk groups defined by \opera{} and by the IPI in
  DLBCL.} Patients with diffuse large B-cell lymphoma and an IPI recorded at
  diagnosis, compared on overall survival, the endpoint the IPI was developed for. \textbf{A}, Kaplan--Meier overall survival by model risk group, with cutoffs fixed on tuning
  data to match the IPI's low-risk observed risk and high-risk specificity
  (\placeholder{cutoffs}). Groups are fixed at baseline, and numbers at risk are shown below the curves. \textbf{B}, The same by IPI category. Shaded bands
  are 95\% bootstrap intervals and the dashed line marks two years. \textbf{C},
  Two-year overall survival per group with 95\% CI, labelled with the number of deaths by two years and the number of patients.
  \textbf{D}, Cross-classification of the two systems. Rows are the IPI
  category and columns the model risk group, so each cell holds the patients
  with that combination, with their two-year survival and 95\% bootstrap
  interval above and the number of patients below. Outlined cells are patients
  placed in the same group by both systems. Cells to the left of the outline
  are patients the model places at lower risk than the IPI, and cells to the
  right at higher risk. Reading along a row shows whether the model separates
  outcomes within one IPI category.}
  \label{fig:risk-groups}
\end{figure}

\FloatBarrier
\operasubsection{What does sharing across diseases and outcomes contribute?}
% Claim heading once results are in, e.g. "Sharing across diseases and outcomes
% each helps, most where a prediction has least data of its own".
\placeholder{Disease pooling / outcome pooling / both improved one-year AUC over
the specialist.} Over the cells evaluable for every model, disease pooling changed one-year AUC by \result{P1_auc}{X}, outcome pooling by \result{P2_auc}{Y} and joint learning by \result{P5_auc}{Z}, with IPA changes of \result{P1_ipa}{A}, \result{P2_ipa}{B} and \result{P5_ipa}{C}. Adding other diseases to an outcome-pooled model changed AUC by \result{P3_auc}{X}, adding other outcomes to a disease-pooled model by \result{P4_auc}{Y}, and the interaction was \result{P6_auc}{Z}, \placeholder{so the two sources were complementary / redundant / interfering} (\cref{fig:pooling}A--B, every cell in \cref{fig:supp-pooling-cells}).

\placeholder{The gain was associated with / not related to the data a
prediction had of its own: shape of the support curve with its interval}
(\cref{fig:pooling}C--D). This curve covers the tasks in which a specialist
could be trained; the sparsest tasks, where it could not, appear as added
coverage below. \placeholder{Trained on X\% of the training patients, the joint
model matched / did not match specialists trained on all of them.} In the subtypes with fewer than 500 training patients,
\opera{}'s one-year AUC differed from per-task XGBoost by
\result{T2_auc_rare}{X}, from pooled XGBoost by \result{T6_auc_rare}{Y} and
from the specialist by \result{P5_auc_rare}{Z} (\cref{tab:panel-effects}). In
the ten subtypes without a dedicated, externally validated index
(\cref{tab:index-availability}), which are assessed today with an index borrowed
from another disease, most often the IPI \citep{IPI1993}, the differences were
\result{T2_auc_noindex}{X}, \result{T6_auc_noindex}{Y} and
\result{P5_auc_noindex}{Z} over \result{T2_cells_noindex}{N} cells, and the IPI
computed from its registry components was compared with \opera{} on the same
patients \placeholder{result} (\cref{fig:score-overview}). \placeholder{Number and
share of cells in which sharing was worse than the specialist, with examples.}
Pooled models also predicted \placeholder{N} cells for which no specialist could
be trained, including every outcome in nodal marginal-zone lymphoma, which had
no patients in the tuning year and so no model of its own, with \placeholder{absolute performance}
(\cref{fig:supp-pooling-cells,fig:supp-horizons}).

\begin{figure}[H]
  \resultfigure{fig_pooling}{Pooling across diseases and outcomes}
    {A: joint versus specialist by disease and domain. B: gain by pooling axis.
    C: gain against training events. D: gain by clinical domain.}
  \caption{\textbf{Effect of pooling across diseases and outcomes on
  discrimination.} \textbf{A}, One-year AUC gain of the joint model over the
  specialist, by disease (rows) and clinical domain (columns), median over the
  domain's outcomes. Diseases run from fewest to most training patients, given beside each name with a log-scale bar, and domains from fewest to most training events.
  \textbf{B}, For each disease, the median gain from pooling
  over diseases, pooling over outcomes and both (joint), on the same rows.
  \textbf{C}, Gain against the number of training events of the target cell,
  with the hierarchical support curve and 95\% credible band. \textbf{D}, Gains
  by clinical domain (median and interquartile range across cells). Positive
  values favour the pooled model, and grey cells could not be estimated.
  Cell-level maps for each pooling axis, their
  interaction and all horizons are in
  \cref{fig:supp-pooling-cells,fig:supp-horizons}.}
  \label{fig:pooling}
\end{figure}

\FloatBarrier
\operasubsection{Does learning what to borrow make sharing work?}
% Claim heading once results are in, e.g. "Letting the model choose its sources
% keeps the gains of sharing and avoids its costs".
The joint model gives every prediction the same shared signal. The routed model
learns, for each patient and target, which diseases and outcomes to draw on
(\cref{fig:source-router}). Routing changed one-year AUC by \result{R1_auc}{X}
and IPA by \result{R1_ipa}{Y} relative to the joint model, and had the higher
AUC in \result{R1_better}{N of M} cells (\cref{fig:routing}A). The gain was not
added capacity: the router without adapters differed from the joint model by
\result{R0_auc}{X}. \placeholder{Nor was it knowing the diagnosis, or having
task-specific capacity: a joint model told the disease, and models restricted to
the patient's own disease or outcome adapter, accounted for X and Y of the
routing gain, so Z came from borrowing from other diseases and outcomes.}
\placeholder{Disease adapters, outcome adapters and both changed AUC by X, Y and
Z} (\cref{fig:routing}B). \placeholder{In cells where the joint model had been
worse than the specialist on the tuning set, routing changed test AUC by X}
(\cref{fig:supp-rescue}), \placeholder{so learning what to borrow recovered
where sharing everything had hurt.} \placeholder{Routing weights concentrated on
the patient's own disease / related diseases} (\cref{fig:supp-router-usage}).

\FloatBarrier
\operasubsection{Can population-scale pretraining extend the shared signal?}
Replacing the encoder pretrained on the cohort with the national encoder, which
had also learned from the rest of the population, changed the joint model's
one-year AUC by \result{E1_auc}{X} when transferred directly and by
\result{E3_auc}{Z} after adaptation to the cohort, and the routed model's by
\result{E2_auc}{Y} and \result{E4_auc}{W} (\cref{fig:encoder-transfer}A).
\placeholder{The incremental effect of adaptation, and whether population
pretraining and routing were additive.} \placeholder{The effect by the share of a
patient's history covered by the national vocabulary, by training events and by
disease} (\cref{fig:encoder-transfer}C--D, \cref{fig:supp-national-where}).
With 10\%, 25\% and 50\% of the training patients, the adapted national encoder
differed from the cohort encoder by \result{N10_auc}{A}, \result{N25_auc}{B} and
\result{N50_auc}{C} (\cref{fig:encoder-transfer}B). \placeholder{The smallest
fraction that reached the full-data cohort-encoder model within the
prespecified margin, or that this was not reached.} As references, pretraining
on the cohort changed the joint model's AUC by \result{P0_auc}{X} over random
initialisation, and fine-tuning by \result{P7_auc}{Y} over a frozen encoder.

\begin{figure}[H]
  \resultfigure{fig_external}{National pretraining and adaptation to the cohort}
    {A: from cohort to national pretraining. B: performance with fewer training
    patients. C--D: the gain by training events and by disease.}
  \caption{\textbf{National pretraining, adaptation to the cohort and
  training-set size.} \textbf{A}, One-year AUC of the joint model (filled) and the routed model
  (outline) initialised from the encoder pretrained on the cohort alone, from
  the national encoder (pretrained on the whole population, the cohort
  included) used directly, and from the national encoder adapted to the cohort
  by continued pretraining on its records (95\% CI,
  averaged over the cells evaluated for all six models). Arrows give the paired
  effect of each step, and the bracket the two together. \textbf{B}, Median
  one-year AUC over the same cells of the joint model from the cohort encoder
  and from the adapted national encoder, trained on 10\%, 25\%, 50\% or all of
  the training patients (nested random subsets with the same patients for every
  outcome). Numbers give the median paired difference at each share; the dashed
  line marks the cohort-encoder model trained on all patients, and the open circle the
  share at which the national model reaches it (interpolated on the log scale).
  \textbf{C}, Paired effect of the adapted national versus the cohort encoder
  for each cell, by the number of training events in the cell (line and band:
  median and interquartile range within bins). \textbf{D}, The same effects by
  disease (diamond and bar: median and interquartile range). Tuning and
  test patients were never subsampled.}
  \label{fig:encoder-transfer}
\end{figure}

\FloatBarrier
\operasubsection{Does one model keep each prediction specific?}
% Claim heading once results are in: "One model reads the same history
% differently for each complication".
Sharing raises an obvious concern: a model trained on every outcome might learn
one general measure of frailty and apply it to all. We decomposed each held-out
prediction of \opera{} with integrated gradients into the contributions of the
recorded events, after checking that the attributions reflect the model:
removing the five most credited events lowered the predicted risk more than
removing five random events for \result{ig_faithful_outcomes}{[n]} of
\result{ig_faithful_total}{[n]} check outcomes, attributions of a model with
randomised weights agreed little with those of the trained model (median rank
agreement \result{ig_rho_random_model}{[x]}), and attributions were stable to
the choice of reference (\result{ig_rho_baseline_source_mean}{[x]}).

Across \result{ig_outcomes}{[n]} outcomes in \result{ig_patients}{[n]} held-out
patients, the timing of care alone (when and how often a patient was seen, with
age and subtype, but not what was recorded) carried a median
\result{ig_var_reference}{[x]} of the differences in predicted risk between
patients. What was recorded carried the rest, and which record mattered
depended on the complication: \placeholder{laboratory results for cytopenias and
organ toxicity, medication for infection and diagnoses for cardiovascular
events and death} (\cref{fig:attribution-main}A, every outcome in
\cref{fig:attribution-atlas}). Records from the last 30 days before treatment
carried \result{ig_var_recent}{[x]}. Outcomes of one clinical domain rested on
the same episodes of a patient's history (median agreement
\result{ig_similarity_within}{[x]}), and outcomes of different domains on
different ones (\result{ig_similarity_between}{[x]})
(\cref{fig:attribution-structure}A): the shared model reads one history
differently for each complication. Of \result{ig_expect_total}{[n]} prespecified
effects of clinical findings, \result{ig_expect_supported}{[n]} were seen and
\result{ig_expect_contradicted}{[n]} went the other way
(\cref{fig:attribution-main}C). \placeholder{Examples, such as high LDH raising
predicted death, low neutrophils severe neutropenia and recent antibacterials
serious infection, and any strong pattern that was not expected.} Molecular
findings recorded as pathology codes moved predictions as expected:
\placeholder{the change in predicted log-odds of treatment failure and death with
a recorded TP53 aberration, MYC rearrangement or unmutated IGHV}
(\cref{fig:attribution-main}C).

Read on the timing of care alone, without retraining, held-out AUC changed by
\result{restrict_timeline}{[x]}, without laboratory results by
\result{restrict_no_labs}{[x]}, without medication by
\result{restrict_no_medication}{[x]}, without diagnoses by
\result{restrict_no_diagnoses}{[x]}, without the region and specialty of care
by \result{restrict_no_region}{[x]} and with only the last year of history by
\result{restrict_last_year}{[x]} (\cref{fig:attribution-main}B, by outcome in
\cref{fig:attribution-structure}B). For a single patient, the same computation
shows which episodes of care move each risk and when they happened
(\cref{fig:attribution-main}D, \cref{fig:attribution-patient}): in the
illustrated patient, a recent pneumonia moves the risk of serious infection, the
findings of the diagnostic work-up the risk of death, and low neutrophils days
before treatment the risk of severe neutropenia.

\begin{figure}[H]
  \resultfigure{fig_attribution}{What OPERA relies on}
    {Differences between patients by source, what each part of the record adds,
    findings by domain, and one patient's history read for three risks.}
  \caption{\textbf{What \opera{} relies on.} Rows in A--C are clinical
  domains. Integrated gradients split each held-out patient's prediction into a
  reference, which keeps the timing of care, age and subtype but replaces every
  recorded event by an average one, and the contribution of each event.
  \textbf{A}, The differences in predicted one-year risk between patients
  carried by the timing of care alone (grey) and by each source of records,
  averaged over the outcomes of each domain. \textbf{B}, Change in held-out AUC
  ($\times$100) when the fitted model reads only part of the record, without
  retraining: only the timing of care, no laboratory results, no medication, no
  diagnoses, no region or specialty of care, or only the last year of history. \textbf{C}, Mean change in
  predicted log-odds among the patients with each finding, averaged over the
  outcomes of each domain and shown as the corresponding odds ratio, with findings grouped by source. Boxed cells contain
  a prespecified expected effect, drawn solid where it was seen and dashed where
  it was not; grey cells had too few patients with the finding to report.
  \textbf{D}, One synthetic patient's history read for three risks, with the
  episodes that move them most labelled above the tracks. Each track is one
  outcome, and each circle is an episode of care at its time before treatment,
  sized by how much it changes the predicted one-year risk (percentage points)
  and coloured by direction. The predicted risk is given at the right, with the
  risk for the same timing of care and average findings. Sienna and indigo mark
  effects throughout (raising risk or losing AUC, and lowering risk), and the
  other colours identify sources of records. Every outcome is shown in
  \cref{fig:attribution-atlas,fig:attribution-structure}, and the patient's
  history read for six complications in \cref{fig:attribution-patient}.}
  \label{fig:attribution-main}
\end{figure}

\FloatBarrier
\operasubsection{How does performance vary over time, across regions and between patients?}
The held-out test set spans patients treated from 2023, a period in which first-line treatment and register coverage changed. \placeholder{One-year AUC and calibration-in-the-large of \opera{} and XGBoost in each calendar year of treatment, and whether the gain was stable from the first to the last year.} \placeholder{The gain did / did not depend on characteristics known at the start of treatment: amount of recorded history, age, sex, region and recorded prognostic index} (\cref{fig:who-benefits,fig:supp-regions}). \placeholder{Results with the known-status estimand, with disease evidence restricted to the prediction date, with size-matched clinical cutoffs, and with equal weight per clinical domain.}

\FloatBarrier
\operasection{Discussion}
Clinical prediction is mostly built one task at a time, and we asked what
happens when it is not. One model, learning across 24 lymphoid cancer subtypes and 69
outcomes, was compared with the tools that would otherwise be built for each
task: gradient-boosted trees tuned for every disease, outcome and horizon,
specialists trained on each prediction alone, and the prognostic indices in
their own populations and on their own endpoints. \placeholder{Principal
findings with uncertainty: how \opera{} compared with XGBoost and the IPI; what
other diseases, other outcomes, learning what to borrow and population
pretraining each added; where the gains were largest; and the most informative
null or harmful result.} These results come from a single evaluation of patients
treated after every modelling decision had been fixed, while first-line
treatment and the registers were changing. The study is not the first to improve
on a prognostic index with machine learning \citep{Biccler2018}, nor the first
to apply sequence models to national records
\citep{shmatko2025delphi,savcisens2024life2vec}, and multi-task learning and
pooling across diseases have been tried before
\citep{caruana1997multitask,Oh2021twostage,harutyunyan2019multitask}. What this
study adds is a controlled test across a whole specialty: where shared signal
comes from, how much each source adds, and whether letting the model learn what
to borrow is what makes sharing work, against the tools built one task at a
time.

The controlled comparisons show where shared signal comes from. Another disease
can supply examples of the same complication; another outcome can teach the
encoder the patient's underlying vulnerability; the two can complement each
other, repeat the same information or conflict. \placeholder{Which pattern
occurred here, and how the gain varied with the data a prediction had of its
own.} Sharing everything is not guaranteed to help: multi-task studies of health
records have reported uneven benefits and negative transfer
\citep{harutyunyan2019multitask,mcdermott2020mtl,ding2019phenotyping}, and in our
earlier study, pooling related aggressive lymphomas matched but did not clearly
improve on a model trained on DLBCL alone \citep{werling2026dlbcl}. At the
other extreme, sharing is the only option: nodal marginal-zone lymphoma had five
training patients and none in the tuning year, so no model of its own could be
fitted or selected, and every prediction for it came from a shared model.
\placeholder{Whether and where sharing everything hurt here, and whether learning
what to borrow recovered it.} The routing controls locate the gain:
the gain did not come from added capacity, from knowing the diagnosis or from
task-specific capacity alone, but from borrowing \placeholder{if the results
hold}. A model that chooses its sources can take what helps a prediction and
leave what does not, which may make sharing easier to apply where its effect
cannot be known in advance. The attributions add that sharing did not blur the
predictions into one measure of frailty: one history was read differently for
each complication. The argument is not specific to haematology, although the
evidence is. Haematology offers many related diseases and outcomes, national
data and established indices to compare against, which makes it a demanding
first test; the same design could be applied wherever related prediction tasks
are each short of events, which in clinical prediction is most of them, common
diseases included, and whether the findings carry over is for those studies to
show.

The broadest source of signal is the population. \placeholder{Whether the
national encoder helped directly, only after adaptation, or not at all, and the
share of training patients with which it matched the cohort encoder.} A gain
need not come only from what the encoder learned from other patients: a
national extract also represents each record at finer resolution, keeping
specific codes, such as molecular findings, that a cohort-sized vocabulary
merges. If population pretraining transfers, a foundation model learned from
population records could become shared infrastructure on which individual
specialties build their prediction models, each adding adaptation to its own
patients and labelled outcomes; a reduction in the labels needed would make this
most valuable for small cohorts, rare diseases and new outcomes. Shared models
also change the practical cost of prediction: one \opera{} model replaced
\result{narrow_fits}{N} specialists \placeholder{and N tabular models, at X of
their combined compute and storage}, and is one pipeline to validate, monitor and
update. The conversion of linked records into model inputs, the reconstruction
of treatment and outcomes, and the training and evaluation workflow can be
reused for the next clinical question, so that a new outcome or population
builds on this foundation rather than starting again.

For patients, a shared model gives a profile of risks at the start of treatment
rather than a single prognostic summary. A patient at low risk of treatment
failure may be at high risk of serious infection, and prophylaxis, monitoring
and dose planning differ accordingly. A profile is useful only if its
probabilities are accurate, which is why calibration is reported alongside
discrimination, and more useful if a clinician can see what each risk rests on:
one patient's explanation can show that a high risk of serious infection follows
a recent pneumonia while a high risk of death follows the findings of the
diagnostic work-up (\cref{fig:attribution-patient}). The comparison with the IPI
connects these predictions to current practice, because the indices already set
treatment and trial eligibility \citep{DLG2025DLBCL,Tilly2022POLARIX}. Better
separation of outcomes than the IPI in its own population, for example by
recognising fit older patients as being at low risk, \placeholder{and the net
benefit across decision thresholds} justifies evaluating whether the additional
information improves supportive-care decisions; it does not justify replacing
IPI-based treatment rules, because the model predicts risk under the care
patients received, not the benefit of changing it \citep{vanGeloven2020}. An
index also applies only where it was recorded or can be computed, whereas the
model gives an estimate for everyone with a record, including subtypes and
complications for which no tool exists. Whether explanations of this kind help
clinicians decide needs testing with them \citep{ghassemi2021explainability}.

The test period differed from the training period in more than calendar time.
First-line treatment changed \placeholder{confirm for Denmark: e.g. polatuzumab
vedotin in DLBCL, targeted agents in CLL, daratumumab combinations in myeloma}
\citep{Tilly2022POLARIX}, the national patient register moved from LPR2 to LPR3
in 2019, and laboratory coverage expanded in eastern Denmark from 2016
\citep{brieghel2025dalycare}. Shifts of this kind are a common cause of failure
for clinical prediction models \citep{finlayson2021shift}. \placeholder{If
performance held, this supports robustness over the observed period}, although
not to future changes or to another health system, and a deployed model needs
monitoring and updating \citep{davis2017drift}, simpler for one model than for
a portfolio. Because the registers are national, one model could serve every
hospital if it works equally well in every region, which it \placeholder{did /
did not} (\cref{fig:supp-regions}). CLL-TIM shows that deployment alongside the
electronic health record is feasible in Denmark \citep{Agius2024clltim}; a
deployed model would need its inputs from hospital systems at the time of the
decision rather than from registers updated with a delay, and its performance
confirmed on data as they are then available.

This study has limitations. Validation was temporal within one health system,
and some test patients contributed earlier history to self-supervised
pretraining, so the evaluation concerns later treatment episodes rather than
unseen individuals; another health system would need to adapt and validate its
own model, which the shared workflow makes possible. Outcomes depend on
ascertainment, follow-up and competing death: laboratory thresholds and
treatment-based proxies are not adjudicated adverse events, laboratory risk sets
condition on monitoring, and some subtypes were coded after treatment started.
Training support varied naturally across tasks, so its relation to the benefit
of sharing is an association, not the effect of removing labels. Models were
trained on patients with known status at each horizon, which tends to
overestimate risk where censoring is common, as for long-horizon laboratory
outcomes; all models share this, but calibration at long horizons should be
read with it in mind. The cutoffs matched to the clinical indices were
estimated on the tuning year and carry uncertainty that intervals with fixed
cutoffs do not include. Each model was
trained once, so intervals omit variation between training runs, and all choices
were made on a tuning year of 1,980 patients, which we addressed by keeping them
few and prespecified. The tabular and sequence pipelines differ in history
windows and representation as well as model class, so their comparison
estimates the value of complete pipelines. The national and cohort encoders
differ in population, register deliveries, vocabulary and training budget, and
the national vocabulary was built from code frequencies that include later
years, so these comparisons concern the encoders as delivered rather than scale
alone. The sequence models did not read quality-registry fields, although
pre-treatment LDH appears in the laboratory histories and molecular findings in
the pathology codes \placeholder{confirm}; the registry-augmented XGBoost shows
what these fields add. Region and specialty of care were inputs, and two regions
contribute richer hospital data, so region partly marks which data a patient
has; performance is reported within each region. For most indices the registers
record the category rather than the score, limiting those comparisons to
stratification and discrimination. Attributions describe the fitted model
relative to a reference, and agreement with clinical expectation shows
plausibility, not causation. Finally, the study does not measure whether
decisions or outcomes improve when the predictions are used. The decision
curves estimate net benefit for one action that changes how closely patients
are followed, not which treatment they receive; a model that changed treatment
would need evidence to the standard of that treatment \citep{blythe2026economic}.
Several of the outcomes, such as serious infection and severe neutropenia after
treatment starts, are what prophylaxis is meant to prevent. The natural next
step is the route taken by CLL-TIM in PreVent-ACaLL \citep{preventacall}: the
model selects patients at high risk, and a randomized trial tests the action in
them.

\placeholder{Conclusion. Draft if the results hold: clinical prediction does not
have to be built one task at a time. A model that shares signal across diseases,
outcomes and the population, and learns what to borrow from where, improved on
the tools built for each task alone and made predictions possible where a
single task had too few events. Haematology was the first test; the approach may
apply wherever
related prediction tasks are each short of data. Its value at the bedside now
needs prospective evaluation.}

\section*{Data availability}
According to Danish legislation, the health data used in this study may not be
shared publicly. Within the approved DALY-CARE protocol, the data may be
accessed on a collaborative basis under a data processing agreement, and all
analyses must be performed within the DALY-CARE data resource at the Danish
National Genome Center \citep{brieghel2025dalycare}. The protocol and the
agreement template are available with the DALY-CARE documentation.
\section*{Code availability}
All code, model configurations, outcome definitions and analysis recipes are
publicly available at \placeholder{repository URL and archival DOI}. The
conversion of the registers to the event format, with its configurations, is
available at \url{https://github.com/FGA-DIKU/ehr2meds} \placeholder{licence}.
Trained model weights were fitted on patient-level data and are not released.
\section*{Patient and public involvement}
\placeholder{Statement on patient and public involvement in the design, conduct or reporting of this research.}
\section*{Acknowledgements}
\placeholder{Funding, collaborators and compute resources.}
\section*{Author contributions}
\placeholder{CRediT statement.}
\section*{Competing interests}
\placeholder{Statement.}
